\scanpage{10}
\mspage{-2.32-}

$G$ on $X_1$ extending the given one on $X$, and there is an $M_1\in\operatorname{Pic}^{G}(X_1)$, ample on $X_1$, which \inserted{\handwritten{extends}} $M^N\handwritten{\otimes}f^*L^{N_0N}$ for some $N$. (This $N$ may be taken as the same $N$ as above). Then we have the diagram:

\diagram{\begin{tikzcd}[ampersand replacement=\&]
X\handwritten{\arrow}[r,"{(I_X,I_Y\circ f)}"]\handwritten{\arrow}[d,"f"']\&X_1\times Y_1\handwritten{\arrow}[d,"{{}_{\unsure{-|{}}}\unsure{p|P}_2}"]\\
Y\handwritten{\arrow}[r,hook,"I_Y"']\&Y_1
\end{tikzcd}}

Now if the result is proven when $X$ and $Y$ are proper, we can apply it to $p_2:X_1\times Y_1\handwritten{\longrightarrow} Y_1$, and the sheaves $M_1\handwritten{\otimes}\underline{\mathrm{o}}_{Y_1}$ on $X_1\times Y_1$ and $L_1$ on $Y_1$. Then

\[
\begin{aligned}
f^{-1}\handwritten{\{}Y^s_{(\unsure{o|0})}(L)\handwritten{\}}&=f^{-1}\handwritten{\{}I_Y^{-1}\handwritten{\{}Y^s_{1,(o)}(L_1)\handwritten{\}}\handwritten{\}}\\
&=(I_X\unsure{{;}|{}_{j}|{,}}I_Y\circ f)^{-1}\handwritten{\{}p_2^{-1}\handwritten{\{}Y^s_{1,(o)}(L_1)\handwritten{\}}\handwritten{\}}\\
&\handwritten{\subset}(I_X\mathpunct{;}I_Y\circ f)^{-1}\handwritten{\{}(X_1\times Y_1)^s_{(o)}(M_1\handwritten{\otimes}L_1^n)\handwritten{\}}\ ;\ \text{if }n\geq n_0\\
&\handwritten{\subset}X^s_{(o)}\handwritten{\bigl[}(I_X,I_Y\circ f)^*(M_1\handwritten{\otimes}L_1^n)\handwritten{\bigr]}\qquad\text{by Prop. 1}\quad,\\
&=X^s_{(o)}\handwritten{\bigl[}M^N\handwritten{\otimes} f^*L^{NN_0+Nn}\handwritten{\bigr]}\\
&=X^s_{(o)}\handwritten{\bigl[}M\handwritten{\otimes} f^*L^{N_0+n}\handwritten{\bigr]}.
\end{aligned}
\]

Now, if $X$ and $Y$ are proper over $k$, we note again that $M\handwritten{\otimes} f^*L^n$ is ample on $X$ for $n\geq n_0$. Therefore, we can apply Theorem 2.1 to the computation of $X^s_{(o)}(M\handwritten{\otimes} f^*L^n)$. But if $x\in X_k$, $\handwritten{\delta}\in\Delta(G)$,

% archive sheet number: 9
